\TieudeT{PHONG TỎA VẬT LÍ}
\subsubsection*{PHẦN I. Thí sinh trả lời câu hỏi từ câu 1 đến câu 18. Mỗi câu hỏi thí sinh chỉ chọn một phương án}
\setcounter{ex}{0}
\begin{ex}%book%Câu 170
Phát biểu nào sau đây đúng?
\choice
{Moment lực tác dụng lên vật là đại lượng vô hướng}
{\True Moment lực đối với một trục quay được đo bằng tích của lực với cánh tay đòn của nó}
{Moment lực là đại lượng đặc trưng cho độ mạnh yếu của lực}
{Đơn vị của moment lực là N/m}
\end{ex}
\begin{ex}%book %[Word to LaTeX 3.2]
Trong hệ SI, đơn vị của moment lực là
\choice
{$N/m$}
{$N$ (Newton)}
{Giây $(s)$}
{\True $N.m$}
\end{ex}
\begin{ex}%book%[1P3Y3-1][Loai1] 
Moment lực tác dụng lên một vật có trục quay cố định là đại lượng 
\choice
{đặc trưng cho độ mạnh yếu của lực}
{đặc trưng cho tác dụng làm quay vật của lực và được đo bằng thương của lực và cánh  tay đòn của nó}
{\True đặc trưng cho tác dụng làm quay vật của lực và được đo bằng tích của lực và cánh với tay đòn của nó}
{luôn có giá trị âm} 
\loigiai{
Moment lực tác dụng lên một vật có trục quay cố định là đại lượng đặc trưng cho tác dụng làm quay vật của lực và được đo bằng tích của lực và cánh tay đòn của nó: $M = Fd$.
} 
\end{ex}
\begin{ex}%book%[1P3Y3-1][Loai1]
Một lực tác dụng vào một vật có trục quay cố định mà \textbf{không} làm cho vật quay khi
\choice
{lực có phương vuông góc với trục quay}
{lực có giá đi qua trọng tâm của vật}
{lực có giá đi qua trọng tâm và có phương vuông góc với trục quay của vật}
{\True lực có giá đi qua trục quay hoặc song song với trục quay}
\loigiai{
Lực có giá đi qua trục quay hoặc song song với trục quay thì không có tác dụng làm vật quay.
}
\end{ex}
\begin{ex}%book%Câu 178
Tác dụng các lực $\vec{F}$ có độ lớn như nhau vào cùng một vị trí trên vật nhưng khác hướng. Trường hợp nào moment của lực $\vec{F}$ có tác dụng làm quay vật quanh $O$ là lớn nhất? 
\begin{center}
\begin{tikzpicture}[draw=Mapcolor]
\draw[fill=gray](0,0)--++(90:0.2)--++(0:3)--++(-90:0.2)--cycle;
\draw[fill=black](0,0.1)circle(0.1cm)node[left]{$O$};
\draw[Mapcolor, very thick,->](3,0.1)--++(90:1.2)node[left,black]{$\vec{F}$};
\draw[fill=white](3,0.1)circle(1pt);
\node[below] at (1.5,-0.25){$a)$};
\begin{scope}[shift={(5,0)},rotate=0]
\draw[fill=gray](0,0)--++(90:0.2)--++(0:3)--++(-90:0.2)--cycle;
\draw[fill=black](0,0.1)circle(0.1cm)node[left]{$O$};
\draw[Mapcolor, very thick,->](3,0.1)--++(0:1.2)node[above,black]{$\vec{F}$};
\draw[fill=white](3,0.1)circle(1pt);
\node[below] at (1.5,-0.25){$b)$};
\end{scope}
\begin{scope}[shift={(10,0)},rotate=0]
\draw[fill=gray](0,0)--++(90:0.2)--++(0:3)--++(-90:0.2)--cycle;
\draw[fill=black](0,0.1)circle(0.1cm)node[left]{$O$};
\draw[Mapcolor, very thick,->](3,0.1)--++(30:1.2)node[above,black]{$\vec{F}$};
\draw[fill=white](3,0.1)circle(1pt);
\node[below] at (1.5,-0.25){$c)$};
\end{scope}
\end{tikzpicture}
\end{center}
\choice
{\True a}
{b}
{c}
{cả ba bằng nhau}
\loigiai{Lớn nhất trường hợp a, nhỏ nhất trường hợp b}
\end{ex}
\begin{ex}%book%Câu 176
\immini[thm]{Khi tác dụng một lực $\vec{F}$ vuông góc với cánh cửa, có độ lớn không đổi vào các vị trí khác nhau như hình. Moment lực gây ra tại vị trí nào là lớn nhất?
\choice
{Điểm $A$}
{Điểm $B$}
{Điểm $C$}
{\True Điểm $D$}}{\begin{tikzpicture}[draw=Mapcolor]
\draw[fill=Mapcolor!30](0,0)--++(90:4)--++(0:3)--++(-90:4)--cycle;
\draw[fill=black](0,0)--++(90:4)--++(0:0.1)--++(-90:4)--cycle;
\draw[fill=black](0,3)circle(0.05cm)node[left]{$A$};
\draw[fill=black](1,2.5)circle(0.05cm)node[left]{$B$};
\draw[fill=black](1.5,1)circle(0.05cm)node[left]{$C$};
\draw[fill=black](2,2)circle(0.05cm)node[left]{$D$};
\node[below] at (1.5,-0.25){Tác dụng lực lên các điểm khác nhau của cửa};
\end{tikzpicture}}
\loigiai{}
\end{ex}
\begin{ex}%book%[1P3Y3-1][Loai1] 
Trường hợp nào sau đây \textbf{không} vận dụng quy tắc moment lực?
\choice
{Một người cầm hòn gạch trên tay}
{Một người dùng xà beng để đẩy một hòn đá} 
{\True Một người dùng tay ép cuốn sách vào tường} 
{Một người cầm càng xe cút kít nâng lên}
\loigiai{
Một người dùng tay ép cuốn sách vào tường không phải là ví dụ của vận dụng quy tắc moment lực.
} 
\end{ex}
\begin{ex}%book%[1P3Y3-1][Loai1] 
Muốn cho vật có trục quay cố định nằm cân bằng thì 
\choice
{tổng moment của các lực phải khác 0}
{tổng moment của các lực phải là một vector có giá đi qua trục quay} 
{tổng moment của các lực phải bằng hằng số}
{\True tổng moment của các lực có xu hướng làm vật quay theo một chiều phải bằng  tổng moment của tất cả các lực có xu hướng làm vật quay theo chiều ngược lại}
\loigiai{
Muốn cho vật có trục quay cố định nằm cằn bằng thì tống moment của các lực có khuynh hướng làm vật quay theo một chiều phải bằng tổng moment của tất cả các lực có khuynh hướng làm vật quay theo chiều ngược lại.
} 
\end{ex}
\begin{ex}%Tailieuchuan
Khi mở hoặc đóng cánh cửa (loại có bản lề), để cánh cửa dễ quay nhất thì ta tác dụng lực vào cạnh
\choice
{gần bản lề, theo hướng vuông góc với mặt phẳng cánh cửa}
{xa bản lề, theo hướng song song với mặt phẳng cánh cửa}
{\True xa bản lề, theo hướng vuông góc với mặt phẳng cánh cửa}
{gần bản lề, theo hướng song song với mặt phẳng cánh cửa}
\loigiai{

}
\end{ex}
\begin{ex}%book%Câu 23
\immini[thm]{Cho biết $F=100\ N,\ OA=100\ cm$. Bỏ qua trọng lượng của thanh. Moment của lực $F$ đối với trục quay $O$ là
\choice
{$50\sqrt{5}\ N.m$}
{$50\sqrt{2}\ N.m$}
{\True $50\sqrt{3}\ N.m$}
{$50\ N.m$}}{\begin{tikzpicture}[scale=1,thick,>=stealth',draw=Mapcolor]
\tkzInit[ymin=-0.5,ymax=4,xmin=-0,xmax=4]\tkzClip
\tkzDefPoints{0/0/m, 3/0/n, 3/-0.1/p, 0/-0.1/q, 0.5/0/o}
\tkzDefShiftPoint[o](60:0.6){g}
\tkzDefShiftPoint[o](60:3){a}
\fill[pattern=north east lines](m)--(n)--(p)--(q);
\tkzDrawSegments(m,n)
\tkzDrawSegments[black, line width=3pt](a,o)
\draw [very thick,Mapcolor,->](a) -- ($(a)!1cm!-60:(o)$)node[above]{$\vec{F}$}coordinate(f);
\tkzDrawPoints[size=3,fill=white](o,a)
\tkzDefPointBy[projection=onto m--n](a) \tkzGetPoint{h}
\tkzDrawSegments[dashed](a,h)
\tkzMarkAngles[size=0.7cm](o,a,h)
\tkzLabelAngles[pos=1.2](o,a,h){$30^\circ$}
\tkzMarkRightAngle[size=0.25](a,h,o)
\tkzLabelPoint[right](a){$A$}
\tkzLabelPoint[below](o){$O$}
\tkzLabelPoint[below](h){$H$}
\end{tikzpicture}}
\haidong{5}	\loigiai{$50\sqrt{3}\ N.m$}
\end{ex}
\begin{ex}%book%Câu 180
\immini[thm]{Một xe cẩu có chiều dài cần trục $\ell=20\ m$ và nghiêng $30^\circ $ so với phương thẳng đứng. Đầu cần trục có treo một thùng hàng nặng $2$ tấn như hình vẽ. Lấy $g=9,8\ m/s^2$. Moment lực do thùng hàng tác dụng lên đầu cần trục đối với trục quay đi qua đầu còn lại của cần trục gắn với thân máy bằng 
\choice
{\True $196000\ N.m$}
{$296000\ N.m$}
{$10000\ N.m$}
{$96000\ N.m$}}{\includegraphics[width=4cm]{"images/img-76"}
}
\haidong{8}	\loigiai{$196000\ N.m$
\begin{itemize}
\item 
\item 
\item 
\item 
\item 
\item 
\item 
\end{itemize}	
}
\end{ex}
\begin{ex}%book%[1P3B3-2][Loai1] 
Một người dùng búa để nhổ một chiếc đinh, khi người đó tác dụng một lực lớn hơn $50\ N$ vào đầu búa thì đinh bắt đầu chuyển động. Biết cánh tay đòn của lực tác dụng của người đó là $20\ cm$ và cánh tay đòn của lực của gỗ tác dụng vào đinh là $2\ cm$. Lực cản lớn nhất của gỗ tác dụng vào đinh bằng
\choice
{\True $500\ N$}
{$350\ N$}
{$200\ N$}
{$400\ N$}
\haidong{6}	\loigiai{
Điều kiện cân bằng: $M_{\text{thuận}} = M_{\text{nghịch}} \Leftrightarrow F_1d_1 = F_2d_2 \Rightarrow F_2=F_1\dfrac{d_1}{d_2}= 50. \dfrac{20}{2} = 500\ N$.
} 
\end{ex}
\begin{ex}%book%Câu 3
\immini[thm]{Một thanh đồng chất có chiều dài $L$, trọng lượng $200\ N$, treo một vật có trọng lượng $450\ N$ vào thanh như hình. Lực $\vec{F}_1$ của thanh tác dụng lên điểm tựa có độ lớn là
\choice
{\True $213\ N$}
{$231\ N$}
{$238\ N$}
{$439\ N$}}{\begin{tikzpicture}[scale=1,thick,>=stealth',draw=Mapcolor]
\tkzDefPoints{0/0/o1, 5/0/o2, 3.475/0.15/c, 2.25/0.15/g, 3.475/-1.5/m, 3.475/-1.25/m'}
\draw(-0.05,0.3)--(-0.05,0)--(5.05,0)--(5.05,0.3)--cycle;
\draw[fill=gray](o1)--++(-60:0.4)--++(180:0.4)--cycle;
\draw[fill=gray](o2)--++(-60:0.4)--++(180:0.4)--cycle;
\tkzDrawPoints[size=3,fill=black](c)
\tkzDrawPoints[size=3,fill=black](g)
\tkzDrawSegments(c,m')
\node at (m) [shape=rectangle,draw=red,pattern=dots,minimum height=0.5cm,minimum width=0.7cm, rounded corners=1pt,label=left:] {};
\draw[dashed,<->](0,0.5)--node[midway,above]{$\dfrac{L}{2}$}(2.25,0.5);
\draw[dashed,<->](2.25,0.5)--node[midway,above]{$\dfrac{L}{4}$}(3.475,0.5);
\draw[->,very thick,Mapcolor](o1)++(90:0.3)--++(90:1)node[right,black]{$\vec{F}_1$};
\draw[->,very thick,Mapcolor](o2)++(90:0.3)--++(90:1.6)node[left,black]{$\vec{F}_2$};
\end{tikzpicture}}
\haidong{5}\end{ex}
\begin{ex}%book%[1P3K3-2][Loai1] 
\immini[thm]{Đặt thanh $AB$ có khối lượng không đáng kể nằm ngang, đầu $A$ gắn vào tường nhờ một bản lề, đầu $B$ nối với tường bằng dây $BC$. Treo vào $B$ một vật có khối lượng $5\ kg$. Cho $AB = 40\ cm$, $AC = 60\ cm$ như hình vẽ. Lấy $g = 10\ m/s^2$. Lực căng $T$ của dây $BC$ có độ lớn là
\choice
{$T = 50\ N$}
{ \True $T = 60\ N$}
{$T = 33,3\ N$}
{$T = 80\ N$}}{\begin{tikzpicture}[scale=1,thick,>=stealth',draw=Mapcolor]
\tkzDefPoints{0/0/c, 0/0.5/m, 0/-4/n, -0.1/-4/p, -0.1/0.5/q}
\tkzDefShiftPoint[c](-90:3){a}
\tkzDefShiftPoint[a](0:2){b}
\tkzDefShiftPoint[b](-90:1){d}
\fill[pattern=north east lines](m)--(n)--(p)--(q);
\tkzDrawSegments(b,c b,d m,n)
\tkzDrawSegments[gray, line width=3pt](a,b)
\tkzDrawPoints[size=3,fill=white](a,b)
\shade[ball color=Mapcolor] (d) circle (1.3ex);
\node[white] at (d) {\normalsize{m}};
\tkzLabelPoint[left](c){C}
\tkzLabelPoint[left](a){A}
\tkzLabelPoint[right](b){B}
\end{tikzpicture} }
\haidong{10}	\loigiai{
\immini{\begin{itemize}
\item Ta có: $BC=\sqrt{AB^2+AC^2}=\sqrt{40^2+60^2}=72\ cm$.
\item Điều kiện cân bằng:\\ $M_P=M_T \Rightarrow T.AH=P.AB \Rightarrow T=P.\dfrac{AB}{AH}=P.\dfrac{BC}{AC}=5.10.\dfrac{72}{60}=60\ N$.
\end{itemize}}{\begin{tikzpicture}[scale=1,thick,>=stealth']
\tkzInit[ymin=-5,ymax=0.5,xmin=-0.5,xmax=3]\tkzClip
\tkzDefPoints{0/0/c, 0/0.5/m, 0/-4/n, -0.1/-4/p, -0.1/0.5/q}
\tkzDefShiftPoint[c](-90:3){a}
\tkzDefShiftPoint[a](0:2){b}
\tkzDefShiftPoint[b](-90:1){d}
\tkzDefPointBy[projection=onto b--c](a)\tkzGetPoint{h}
\tkzDefPointBy[homothety=center b ratio 0.25](c)\tkzGetPoint{t}
\tkzDefPointBy[projection=onto b--d](t)\tkzGetPoint{t'}
\tkzDefPointBy[translation=from t' to b](d)\tkzGetPoint{p'}
\fill[pattern=north east lines](m)--(n)--(p)--(q);
\tkzDrawSegments(b,c b,d m,n)
\tkzDrawSegments[dashed,thin](h,a)
\tkzMarkRightAngle(c,h,a)
\tkzDrawSegments[gray, line width=3pt](a,b)
\tkzDrawPoints[size=3,fill=white](a,b)
\shade[ball color=black] (d) circle (1.3ex);
\tkzDrawSegments[blue,very thick,->](b,t d,p')
\node[white] at (d) {\normalsize{m}};
\tkzLabelPoint[left](c){C}
\tkzLabelPoint[left](a){A}
\tkzLabelPoint[right](b){B}
\tkzLabelPoint[above,xshift=5pt](h){H}
\tkzLabelPoint[right](t){$\vec{T}$}
\tkzLabelPoint[right](p'){$\vec{P}$}
\end{tikzpicture} }
} 
\end{ex}
\begin{ex}%Tailieuchuan
\immini[thm]{Cho hệ như hình vẽ. Thanh $AC$ đồng chất, tiết diện đều có trọng lượng $3\ N$. Vật treo tại A có trọng lượng $8\ N$. Lấy $g=10 \ m/s^2$. Để hệ cân bằng nằm ngang, lực $F$ đặt tại $B$ phải có độ lớn
\choice
{$2,5 \ N$}
{$7,5 \ N$}
{$10\ N$}
{\True $5\ N$}}{\begin{tikzpicture}[scale=1,thick,>=stealth',draw=Mapcolor]
\tkzDefPoints{0/0/o1, 5/0/o2, 3.475/0.15/c, 2.25/0.15/g,0/-1.5/m, 3.475/-1.25/m'}
\draw(-0.05,0.3)--(-0.05,0)--(4,0)--(4,0.3)--cycle;



\node at (m) [shape=rectangle,draw=red,pattern=dots,minimum height=0.5cm,minimum width=0.7cm, rounded corners=1pt,label=left:] {};

\node at (0.5,0) {};
\node at (1,0) {};
\node at (1.5,0) {};
\node at (2,0) {};
\node at (2.5,0) {};
\node at (3,0) {};
\node at (3.5,0) {};
\draw (0.5,0.3) -- (0.5,0);

\draw (1,0.3) -- (1,0);

\draw (1.5,0.3) -- (1.5,0);
\draw (2,0.3) -- (2,0);
\draw (2.5,0) -- (2.5,0.3);
\draw (3,0) -- (3,0.3);

\draw (3.5,0) -- (3.5,0.3);
\draw (0,0) -- (0,-1.22);
\node at (-0.1,0.5) {A};
\node at (1,0.5) {O};
\node at (2,0.5) {B};
\node at (4,0.5) {C};
\draw[fill=gray!30] (1,0) -- (0.9,-0.2) -- (1.1,-0.2) -- (1,0);
\draw[->,Mapcolor] (2,0) -- (2,-0.8);
\node at (2.3,-0.9) {$\vec F$};
\end{tikzpicture}}
\loigiai{}
\end{ex}
\begin{ex}%Tailieuchuan
\immini[thm]{Một người dùng một lực $F=50 \ N$ tác dụng vuông góc vào đầu cờ lê với khoảng cách từ điểm đặt lực $F$ tới trục quay qua $O$ như hình vẽ bên dưới là $0,3 \ m$ để xoay ốc. Moment của lực F tác dụng lên cờ lê là
\choice
{$17\ N.m$}
{$12\ N.m$}
{\True $15\ N.m$}
{$10\ N.m$}
}
{\includegraphics[width=3.5cm]{images/2025_04_24_a80008410be527bbd242g-13}}
\loigiai{}
\end{ex}

\begin{ex}%book%[1P3K3-2][Loai1] 
\immini[thm]{Một thanh gỗ dài $1,8\ m$ nặng $30\ kg$, một đầu được gắn vào trần nhà nhờ một bản lề, đầu còn lại được buộc vào một sợi dây và gắn vào trần nhà sao cho phương của sợi dây thẳng đứng và giữ cho tấm gỗ nằm nghiêng hợp với trần nhà nằm ngang một góc $45^\circ$. Biết trọng tâm của thanh gỗ cách đầu gắn sợi dây $60\ cm$. Lấy $g=10\ m/s^2$. Lực căng của sợi dây là
\choice
{$240\ N$}
{\True $200\ N$}
{$300\ N$}
{$100\ N$}}{\begin{tikzpicture}[scale=1,thick,>=stealth',draw=Mapcolor]
\tkzInit[ymin=-2.5,ymax=0.5,xmin=-0,xmax=4]\tkzClip
\tkzDefPoints{0/0/m, 3.5/0/n, 3.5/0.1/p, 0/0.1/q, 3/0/o}
\tkzDefShiftPoint[o](225:1){g}
\tkzDefShiftPoint[o](225:3){a}
\tkzDefPointBy[projection=onto m--n](a)\tkzGetPoint{h}
\fill[pattern=north east lines](m)--(n)--(p)--(q);
\tkzDrawSegments(m,n a,h)
\tkzDrawSegments[Mapcolor, line width=3pt](a,o)
\tkzDrawPoints[size=3,fill=white](o,a)
\tkzMarkAngles[size=0.4cm](h,o,a)
\tkzLabelAngles[pos=-0.8](h,o,a){$45^\circ$}
\end{tikzpicture}} 
\haidong{8}		\loigiai{
\immini{Điều kiện cân bằng $M_{\text{thuận}} = M_{\text{nghịch}}\\
\Rightarrow T.d' = P.d
\Rightarrow T. OA.\cos 45^\circ = P. OG. \cos 45^\circ
\Rightarrow T.1,8 = 30.10.1,2
\Rightarrow T = 200\ N$}{\begin{tikzpicture}[scale=1,thick,>=stealth']
\tkzInit[ymin=-3.5,ymax=0.5,xmin=-0,xmax=4]\tkzClip
\tkzDefPoints{0/0/m, 3.5/0/n, 3.5/0.1/p, 0/0.1/q, 3/0/o}
\tkzDefShiftPoint[o](225:1){g}
\tkzDefShiftPoint[o](225:3){a}
\tkzDefPointBy[projection=onto m--n](a)\tkzGetPoint{h}
\fill[pattern=north east lines](m)--(n)--(p)--(q);
\tkzDrawSegments(m,n a,h)
\tkzDrawSegments[gray, line width=3pt](a,o)
\tkzMarkAngles[size=0.4cm](h,o,a)
\tkzLabelAngles[pos=-0.8](h,o,a){$45^\circ$}
\draw[very thick,blue,->](g)--++(-90:1.8)node[right]{$\vec{P}$};
\draw[very thick,blue,->](a)--++(90:0.6)node[left]{$\vec{T}$};
\tkzDrawPoints[size=3,fill=white](o,a,g)
\tkzLabelPoint[above](o){O}
\tkzLabelPoint[left](a){A}
\tkzLabelPoint[below right](g){G}
\end{tikzpicture}}
} 
%<MyLT1>
\end{ex}
\begin{ex}%book%[1P3K3-2][Loai1]
Một thanh sắt dài, đồng chất, tiết diện đều, được đặt trên bàn sao cho $\dfrac{1}{4}$ chiều dài của nó nhô ra khỏi bàn. Tại đầu nhô ra, người ta đặt một lực $ \vec{F} $ hướng thẳng đứng xuống dưới. Khi lực đạt tới giá trị lớn hơn $40\ N$ thì đầu kia của thanh sắt bắt đầu bênh lên. Lấy $ g=10\ m/s^2 $. Khối lượng của thanh là
\choice
{\True $4\ kg$}
{$5\ kg$}
{$6\ kg$}
{$7\ kg$}
\haidong{10}	\loigiai{
\immini{\begin{itemize}\item Xét trục quay là điểm tiếp xúc O giữa mép bàn và thanh sắt. 
\item Thanh sắt đồng chất, tiết diện đều: 
$ AG=GB\; \Rightarrow GO=OB=\dfrac 1 4AB $
\item Khi đầu kia của thanh sắt bắt đầu bênh lên, ta có: 
$ M_F=M_P\Rightarrow F.OB=P.OG=mg.OG $ 
$ \Rightarrow m=\dfrac {F.OB}{g.OG}=4\ kg.$ 
\end{itemize}}{\begin{tikzpicture}[scale=1,thick,>=stealth']
\draw[pattern=north east lines](0,0)node[above=3pt]{\normalsize{A}} rectangle (4,0.1) node[above=0pt]{\normalsize{B}};
\draw[fill=gray](-0.5,-0.04) rectangle (3,-0.24);
\draw[fill=gray](2.5,-0.24) rectangle (2.7,-1);
\node[above=3pt] at (2,0){\normalsize{G}};
\node[above=3pt,above] at (3,0){\normalsize{O}};
\coordinate (o) at (3,0.05);
\draw[fill=black] (o) node[left]{}circle(0.05);
\draw[->] (2,0)--++(270:1) coordinate(m) node[left]{\normalsize{\bth{\vec{P}}}};
\draw[->] (4,0)--++(270:1) coordinate(m2) node[right]{\normalsize{\bth{\vec{F}}}};
\end{tikzpicture} }}
\end{ex}










\subsubsection*{PHẦN 2. Thí sinh trả lời từ câu 1 đến câu 4. Trong mỗi ý a), b), c), d) ở mỗi câu, thí sinh chọn đúng hoặc sai.}
\setcounter{ex}{0}
\begin{ex}%Tailieuchuan
Một vật chịu tác dụng của lực có xu hướng làm cho vật quay quanh một trục cố định.
\choiceTF[t]
{\True Moment lực tỉ lệ thuận với độ lớn của lực tác dụng} 
{\True Moment lực sẽ luôn bằng $0$ nếu lực tác dụng song song với trục quay} 
{Moment lực không phụ thuộc vào hướng của lực tác dụng} 
{\True Khi khoảng cách từ điểm tác dụng lực đến trục quay càng lớn thì moment lực càng lớn} 
\loigiai{ 
\begin{itemchoice}
\itemch
\itemch
\itemch
\itemch
\end{itemchoice}
}
\end{ex}
\begin{ex}%Tailieuchuan
Xét một vật rắn quay quanh trục:
\choiceTF[t]
{\True Khi Một vật rắn quay quanh một trục thì tổng moment lực tác dụng lên vật có giá trị bằng 0} 
{\True Muốn cho một vật có trục quay cố định nằm cân bằng thì tổng moment của các lực có khuynh hướng làm vật quay theo một chiều phải bằng tổng moment của các lực có khuynh hướng làm vật quay theo chiều ngược lại} 
{Muốn cho một vật có trục quay cố định nằm cân bằng thì tổng moment của các lực phải là một vecto có giá đi qua trục quay} 
{Hệ hai lực cân bằng và ba lực cân bằng có chung tính chất cùng giá và cùng độ lớn} 
\loigiai{ 
\begin{itemchoice}
\itemch
\itemch
\itemch
\itemch
\end{itemchoice}
}
\end{ex}
\begin{ex}%Tailieuchuan
Một vật có khối lượng $5\ kg$ được buộc vào đầu một chiếc gậy dài $90\ cm$. Một người quẩy lên trên vai chiếc gậy đó biết vật cách vai một khoảng là $60\ cm$. Đầu còn lại của chiếc gậy được giữ bằng tay theo phương vuông góc với gậy. Bỏ qua trọng lượng của gậy, lấy $g=10 \ m/s^2$.
\choiceTF[t]
{\True Lực do vật tác dụng lên gậy là $50\ N$} 
{Moment lực do vật tác dụng lên gậy có độ lớn $300\ N.m$} 
{Để gậy được giữa cân bằng thì lực do vật tác dụng lên gậy phải bằng lực do tay tác dụng lên gậy} 
{Lực giữ của tay bằng $25\ N$} 
\loigiai{ 
\begin{itemchoice}
\itemch
\itemch
\itemch
\itemch
\end{itemchoice}
}
\end{ex}
\begin{ex}%book
Thanh $AB$ đồng chất, tiết diện đều có trọng lượng $P=10\ N$. Người ta treo các vật có trọng lượng $P_{1}=20\ N$ và $P_{2}=30\ N$ lần lượt tại $A$, $B$ và đặt giá đỡ tại $O$ để thanh cân bằng. Biết $AB=1,2\ m$.
\choiceTF[t]
{\True Trọng lực $\vec{P}_1$ của vật và trọng lực $\overrightarrow{{P}}$ của thanh làm thanh quay ngược chiều kim đồng hồ quanh trục O}
{\True Trọng lực $\overrightarrow{P_2}$ làm thanh quay quanh cùng chiều kim đồng hồ quanh trục O}
{Momen của phản lực có giá đi qua $O$ là $10\ N.m$}
{Giá đỡ cách đầu A một khoảng bằng $0,6\ m$ thì hệ cân bằng}
\loigiai{\begin{itemchoice}
\itemch \immini{a) Đúng.
Lực tác dụng lên thanh AB như hình bên

Từ hình vẽ ta thấy: Trọng lực $\vec{P}_{1}$ của vật và trọng lực

$\overrightarrow{P}$ của thanh làm thanh quay ngược chiều kim đồng hồ quanh trục O}{\begin{tikzpicture}[scale=1,thick,>=stealth']
\draw(0,0) rectangle (4,0.2);
\node [above=3pt] at (0,0){\normalsize{A}};
\node [above right=3pt] at (2.5,0){\normalsize{O}};
\node [above=3pt] at (4,0){\normalsize{B}};
\node [above=3pt] at (2,0){\normalsize{G}};
\draw[pattern=north west lines] (2.5,0)--++(-60:0.5)--++(180:0.5)--cycle;
\draw[->] (0,0)--++(270:0.8) coordinate(m) node[right]{\normalsize{\bth{\vec{P}_1}}};
\draw[->] (2,0)--++(270:0.6) coordinate(m) node[left]{\normalsize{\bth{\vec{P}}}};
\draw[->] (4,0)--++(270:1.6) coordinate(m) node[right]{\normalsize{\bth{\vec{P}_2}}};
\draw[->] (2.5,0)--++(90:3) coordinate(m2) node[right]{\normalsize{\bth{\vec{N}}}};
\end{tikzpicture}}

\itemch Đúng.
Từ hình vẽ ta thấy: Trọng lực $\vec{P}_{2}$ làm thanh quay quanh cùng chiều kim đồng hồ quanh trục O .

\itemch Sai. Phản lực $\vec{N}$ có giá đi qua O nên $M_{\bar{N}}=0$

\itemch Sai. Gọi $M_{\vec{P}_{1}} ; M_{\vec{P}_{2}} ; M_{\vec{P}}$ lần lượt là momen của $\vec{P}_{1}, \vec{P}_{2}$ và $\vec{P}$ đối với trục quay O

Khi hệ cân bằng, áp dụng quy tắc momen, ta có:

$M_{\vec{P}_{1}}+M_{\vec{P}}=M_{\vec{P}_{2}} \Rightarrow P_{1} . O A+P . O I=P_{2} . O B$

Mặt khác: $O I=O A-\dfrac{A B}{2}$ và $O B=A B-O A$

$\Rightarrow P_{1} . O A+P\left(O A-\dfrac{A B}{2}\right)=P_{2} .(A B-O A)$

$\Rightarrow 20 . O A+10 .(O A-0,6)=30 .(1,2-O A) \Rightarrow O A=0,7\ m$

\end{itemchoice}		
}
\end{ex}























\subsubsection*{PHẦN 3. Thí sinh trả lời từ câu 1 đến câu 6.}
\setcounter{ex}{0} 
\begin{ex}%Tailieuchuan
\immini[thm]{Một chiếc xe đẩy chuyển vật liệu có cấu tạo như hình vẽ. Tổng khối lượng vật liệu và xe là $200\ kg$. Lấy $g=9,8\ m/s^2$. Lực nâng đặt vào tay cầm bằng bao nhiêu niutơn để giữ xe thăng bằng?}{
\includegraphics[width=4cm]{images/2025_04_24_a80008410be527bbd242g-11}}
\shortans[oly]{840}
\loigiai{}
\end{ex}

\begin{ex}%Tailieuchuan
Moment lực của một lực $5\ N$ đối với trục quay bằng $40\ N.m$ thì cánh tay đòn dài bao nhiêu mét?
\shortans[oly]{8}
\loigiai{}
\end{ex}

\begin{ex}%Tailieuchuan
\immini[thm]{Người ta tác dụng lực $\vec{F}$ có độ lớn $100\ N$ lên tay quay để xoay chiếc cối xay như hình. Cho rằng $\vec{F}$ có phương tiếp tuyến với bề mặt cối xay, khoảng cách từ tay quay đến tâm quay là $d=50\ cm$. Moment của lực $\vec{F}$ đối với trục quay qua tâm cối xay bằng bao nhiêu $N.m$?}{
\includegraphics[width=4cm]{images/2025_04_24_a80008410be527bbd242g-12(1)}}
\shortans[oly]{50}
\loigiai{}
\end{ex}

\begin{ex}%Tailieuchuan
Thanh $AC$ đồng chất có trọng lượng $4\ N$, chiều dài $8\ cm$. Biết quả cân có trọng lượng $P_1=10\ N$ treo vào đầu $A$, quả cân có trọng lượng $P_2$ treo vào đầu $C$. Trục quay cách $A$ một đoạn $2\ cm$. Hệ cân bằng thì $P_2$ có độ lớn bằng bao nhiêu niutơn?
\shortans[oly]{2}
\loigiai{}
\end{ex}

\begin{ex}%Tailieuchuan
\immini[thm]{Một người nâng tấm ván $AB$ có khối lượng $m$ (phân bố đều theo chiều dài) với lực $F_1=200\ N$ theo phương vuông góc với tấm ván như hình vẽ. Nếu thay lực $F_1$ bằng lực $F_2$ theo phương thẳng đứng hướng lên thì $F_2$ phải bằng bao nhiêu niutơn thì thanh vẫn cân bằng (làm tròn kết quả đến chữ số hàng đơn vị)?}{
\begin{tikzpicture}[scale=1,draw=Mapcolor, very thick,>=stealth',draw=Mapcolor]
\tkzInit[ymin=-0.5,ymax=4,xmin=-0,xmax=4]\tkzClip
\tkzDefPoints{0/0/m, 3.5/0/n, 3.5/-0.1/p, 0/-0.1/q, 0.5/0/o}
\tkzDefShiftPoint[o](60:0.6){g}
\tkzDefShiftPoint[o](30:3){a}
\fill[pattern=north east lines](m)--(n)--(p)--(q);
\tkzDrawSegments(m,n)
\tkzDrawSegments[gray, line width=3pt](a,o)
\tkzMarkAngles[size=0.7cm](n,o,a)
\tkzLabelAngles[pos=1.0](n,o,a){$30^\circ$}
\draw [very thick,Mapcolor,->](a) -- ($(a)!1cm!-90:(o)$)node[above]{$\vec{F}$}coordinate(f);
\tkzDrawPoints[size=3,fill=white](o,a)
\tkzLabelPoint[above left](o){$A$}
\tkzLabelPoint[right](a){$B$}
\draw[very thick,->,Mapcolor] (2,0.9) -- (2,0.2);
\node at (2.5,0.5) {$\vec{P}$};
\node at (1.2,-0.3) {$d_P$};
\end{tikzpicture}}
\shortans[oly]{231}
\loigiai{
}
\end{ex}

\begin{ex}%Tailieuchuan
Moment lực của một lực đối với trục quay là bao nhiêu $N.m$ nếu độ lớn của lực là $5,5 \ N$ và cánh tay đòn là $4\ m$?
\shortans[oly]{22}
\loigiai{

}
\end{ex}


\newpage 
\TieudeT{PHONG TỎA VẬT LÍ}
\subsubsection*{PHẦN I. Thí sinh trả lời câu hỏi từ câu 1 đến câu 18. Mỗi câu hỏi thí sinh chỉ chọn một phương án}
\setcounter{ex}{0}
\begin{ex}%book%[1P3Y3-1][Loai1] 
Moment lực được xác định bằng công thức
\choice
{$P = mg$}
{$M = \dfrac{F}{d}$}
{ \True $M = Fd$}
{$F = ma$}
\loigiai{
Moment lực được xác định bằng công thức: $M = Fd$.	 
} 
\end{ex}
\begin{ex}%book%[1P3Y3-1][Loai1]
Moment lực tác dụng lên một vật là đại lượng
\choice
{để xác định cánh tay đòn}
{\True đặc trưng cho tác dụng làm quay của lực}
{để xác định độ lớn của lực tác dụng}
{luôn có giá trị dương}
\loigiai{ 
Moment lực tác dụng lên một vật là đại lượng  đặc trưng cho tác dụng làm quay của lực.
}
\end{ex}
\begin{ex}%book%[1P3Y3-1][Loai1]
Phát biểu nào sau đây đúng với quy tắc moment?
\choice
{\True Muốn cho một vật có trục quay cố định nằm cân bằng thì tổng moment của các lực có khuynh hướng làm vật quay theo một chiều phải bằng tổng moment của các lực có khuynh hướng làm vật quay theo chiều ngược lại}
{Muốn cho một vật có trục quay cố định nằm cân bằng thì tổng moment của các lực phải bằng hằng số}
{Muốn cho một vật có trục quay cố định nằm cân bằng thì tổng moment của các lực phải khác 0}
{Muốn cho một vật có trục quay cố định nằm cân bằng thì tổng moment của các lực phải là một vector có giá đi qua trục quay}
\loigiai{
Muốn cho một vật có trục quay cố định nằm cân bằng thì tổng moment của các lực có khuynh hướng làm vật quay theo một chiều phải bằng tổng moment của các lực có khuynh hướng làm vật quay theo chiều ngược lại  đúng với quy tắc moment.
}
\end{ex}
\begin{ex}%book%[1P3Y3-1][Loai1]
Trong trường hợp nào sau đây, lực tác dụng vào vật làm cho vật quay quanh một trục?
\choice
{Lực $\vec F_1 $ có phương song song với trục quay}
{Lực $\vec F_2 $ có phương đi qua trục quay}
{\True Lực $\vec F_3 $ có phương không đi qua trục quay, nằm trong mặt phẳng vuông góc với trục quay}
{Tất cả các đáp án trên}
\loigiai{
Chỉ có lực $\vec F_3 $ có phương không đi qua trục quay, nằm trong mặt phẳng vuông góc với trục quay  mới làm cho vật quay.
}
\end{ex}
\begin{ex}%book %[Word to LaTeX 3.2]
Tác dụng làm quay vật của một một lực \textbf{không} phụ thuộc vào
\choice
{cánh tay của đòn lực}
{độ lớn của lực}
{vị trí của trục quay}
{\True điểm đặt của lực}
\loigiai{}
\end{ex}
\begin{ex}%book%[1P3Y3-1][Loai1] 
Ở trường hợp nào sau đây, lực có tác dụng làm cho vật rắn quay quanh trục? 
\choice
{Lực có giá song song với trục quay}
{Lực có giá nằm trong mặt phẳng vuông góc với trục quay và cắt trục quay}
{ \True Lực có giá nằm trong mặt phẳng vuông góc với trục quay và không cắt trục quay}
{Lực có giá cắt trục quay}
\loigiai{
Lực có giá nằm trong mặt phẳng vuông góc với trục quay và không cắt trục quay thì lực có tác dụng làm cho vật rắn quay quanh trục.
} 
\end{ex}
\begin{ex}%book 
Trường hợp nào sau đây lực tác dụng vào một vật có trục quay cố định mà \textbf{không} làm cho vật quay? 
\choice
{Lực có giá nằm trong mặt phẳng vuông góc với trục quay và không cắt trục quay}
{\True Lực có giá song song với trục quay}
{Tổng đại số mômen các lực tác dụng lên vật khác 0}	
{Mômen của lực tác dụng lên vật theo chiều kim đồng hồ lớn hơn hoặc nhỏ hơn mômen của lực tác dụng lên vật theo chiều ngược lại}
\end{ex} 
\begin{ex}%Tailieuchuan
Đòn bẩy là ứng dụng của quy tắc
\choice
{mặt phẳng nghiêng}
{quán tính}
{\True moment lực}
{đòn gánh}
\loigiai{}
\end{ex}
\begin{ex}%Tailieuchuan
Một vật đang quay quanh một trục cố định với tốc độ góc không đổi. Nếu bỗng nhiên tất cả moment lực tác dụng lên vật mất đi thì
\choice
{vật quay chậm dần rồi dừng lại}
{vật quay nhanh dần đều}
{vật lập tức dừng lại}
{\True vật tiếp tục quay đều}
\loigiai{}
\end{ex}
\begin{ex}%Tailieuchuan
Một tấm ván nặng $270\ N$ được bắc qua một con mương. Trọng tâm của tấm ván cách điểm tựa trái $0,80 \ m$ và cách điểm tựa phải là $1,60 \ m$. Lực mà tấm ván tác dụng lên điểm tựa bên trái là
\choice
{\True $180 \ N$}
{$90\ N$}
{$160\ N$}
{$80\ N$}
\loigiai{}
\end{ex}
\begin{ex}%Tailieuchuan
Một người nông dân dùng quang gánh, gánh hai thúng, thúng gạo và thúng ngô. Thúng gạo nặng $30\ kg$, đòn gánh có chiều dài $1,5 \ m$. Vai người nông dân đặt cách gánh gạo $1\ m$. Thúng ngô có khối lượng là
\choice
{$45\ kg$}
{$15\ kg$}
{\True $60\ kg$}
{$30\ kg$}
\loigiai{}
\end{ex}
\begin{ex}%book%Câu 22
\immini[thm]{Dưới tác dụng của lực $\vec{F}$, thanh $AB$ có độ dài $5\ cm$ có thể quay quanh điểm A. Cánh tay đòn của lực $\vec{F}$ trong trường hợp này là 
\choice
{$3\ cm$}
{$4\ cm$}
{\True $5\ cm$}
{$6\ cm$}}{\begin{tikzpicture}[scale=1,thick,>=stealth',draw=Mapcolor]
\tkzInit[ymin=-0.5,ymax=4,xmin=-0,xmax=4]\tkzClip
\tkzDefPoints{0/0/m, 3.5/0/n, 3.5/-0.1/p, 0/-0.1/q, 0.5/0/o}
\tkzDefShiftPoint[o](60:0.6){g}
\tkzDefShiftPoint[o](30:3){a}
\fill[pattern=north east lines](m)--(n)--(p)--(q);
\tkzDrawSegments(m,n)
\tkzDrawSegments[gray, line width=3pt](a,o)
\tkzMarkAngles[size=0.7cm](n,o,a)
\tkzLabelAngles[pos=1.0](n,o,a){$30^\circ$}
\draw [very thick,Mapcolor,->](a) -- ($(a)!1cm!-90:(o)$)node[above]{$\vec{F}$}coordinate(f);
\tkzMarkRightAngle(f,a,o)
\tkzDrawPoints[size=3,fill=white](o,a)
\tkzLabelPoint[above left](o){$A$}
\tkzLabelPoint[right](a){$B$}
\end{tikzpicture}}
\haidong{5}
\loigiai{$AB=5\ cm.$}
\end{ex}
\begin{ex}%book%Câu 179
\immini[thm]{Biết mỏ lết có chiều dài $25\ cm$ và khoảng cách từ điểm đặt của lực đến bu lông vào cỡ $21\ cm$. Moment do lực $\vec{F}$ có độ lớn $10\ N$ tác dụng vuông góc lên mỏ lết để làm xoay bu lông bằng
\choice
{\True $2,1\ N.m$}
{$1,1\ N.m$}
{$3,3\ N.m$}
{$5,2\ N.m$}}{\begin{tikzpicture}[scale=1,thick,>=stealth']
\tkzInit[ymin=-0.5,ymax=2.5,xmin=-1,xmax=5]\tkzClip
\tkzDefPoints{0/1/o, 2.7/1/a}
\node[inner sep=0pt,rotate=00] (russell) at (-0.3,1.02){\includegraphics[width=0.58cm]{images/ecu.pdf}};
\node[inner sep=0pt,rotate=20] (russell) at (1.7,1){\includegraphics[width=5cm]{images/cole.pdf}};
\tkzDrawPoints[size=2,fill=black](o,a)
\tkzDrawSegments[dashed,blue](o,a)
\draw[very thick,blue,->](a)--++(-90:1.2)coordinate(f)node[right,black]{$\vec{F}$};
\tkzLabelPoint[left=15pt](o){$O$}
\tkzDrawArc[R,arrows=<-,color=black](o,1.2cm)(30,140)
\tkzDefPointBy[projection=onto a--f](o)\tkzGetPoint{h}
\tkzDrawSegments[dashed,thin](o,h)
\tkzLabelSegment(o,h){$d$}
\end{tikzpicture}}
\haidong{5}	\loigiai{$M=F.d=2,1\ N.m$}
\end{ex}
\begin{ex}%book%Câu 175
Một bu lông nối khung chính và khung sau của xe đạp leo núi cần moment lực $15\ N.m$ để siết chặt. Nếu bạn có khả năng tác dụng lực $40\ N$ lên cờ lê theo một hướng bất kì thì chiều dài tối thiểu của cờ lê để tạo ra moment lực cần thiết là
\choice
{\True $0,375\ m$}
{$0,325\ m$}
{$\dfrac{8}{3}\ m$}
{$600\ m$}
\haidong{8}\loigiai{Chiều dà tối thiểu của cờ lê: $d=\dfrac{M}{F}=0,375\ m$}
\end{ex}
\begin{ex}%book%[1P3B3-2][Loai1] 
Một cánh cửa chịu tác dụng của một lực có moment $M_1 = 60\ N.m$ đối với trục quay đi qua các bản lề. Lực $F_2$ tác dụng vào cửa có moment quay theo chiều ngược lại và có cánh tay đòn $d_2 = 1,5\ m$. Lực $F_2$ có độ lớn bằng bao nhiêu thì cửa không quay?
\choice
{$60\ N$}
{\True $40\ N$}
{$90\ N$}
{$120\ N$}
\haidong{7}	\loigiai{
Để cánh cửa không quay thì: $M_1 = F_2d_2 \Leftrightarrow 60 = 1,5.F_2 \Rightarrow F_2 = 40\ N$.
} 
\end{ex}
\begin{ex}%book%[1P3K3-2][Loai1] 
\immini[thm]{Một người nâng một tấm gỗ dài $1,5\ m$, nặng $30\ kg$ và giữ cho nó hợp với mặt đất nằm ngang một góc $60^\circ$. Biết trọng tâm của tấm gỗ cách đầu mà người đó nâng $120\ cm$, lực nâng vuông góc với tấm gỗ. Lấy $g=10\ m/s^2$. Lực nâng của người đó là
\choice
{\True $30\ N$}
{$240\ N$}
{$51,96\ N$}
{$300\ N$}}{\begin{tikzpicture}[scale=1,thick,>=stealth',draw=Mapcolor]
\tkzInit[ymin=-0.5,ymax=4,xmin=-0,xmax=4]\tkzClip
\tkzDefPoints{0/0/m, 3.5/0/n, 3.5/-0.1/p, 0/-0.1/q, 0.5/0/o}
\tkzDefShiftPoint[o](60:0.6){g}
\tkzDefShiftPoint[o](60:3){a}
\fill[pattern=north east lines](m)--(n)--(p)--(q);
\tkzDrawSegments(m,n)
\tkzDrawSegments[gray, line width=3pt](a,o)
\tkzMarkAngles[size=0.4cm](n,o,a)
\tkzLabelAngles[pos=0.6](n,o,a){$\alpha$}
\draw [very thick,Mapcolor,->](a) -- ($(a)!1cm!-90:(o)$)node[above]{$\vec{F}$}coordinate(f);
\tkzMarkRightAngle(f,a,o)
\tkzDrawPoints[size=3,fill=white](o,a)
\end{tikzpicture}}
\haidong{8}	\loigiai{
\immini{Điều kiện cân bằng: $M_{\text{thuận}} = M_{\text{nghịch}}\\
\Rightarrow P.d = F. OA \Leftrightarrow mg. OG. \cos 60^\circ = F. OA
\Rightarrow 30.10.30.0,5 = F.150
\Rightarrow F = 30\ N$}{\begin{tikzpicture}[scale=1,thick,>=stealth']
\tkzInit[ymin=-2,ymax=4,xmin=-0,xmax=4]\tkzClip
\tkzDefPoints{0/0/m, 3.5/0/n, 3.5/-0.1/p, 0/-0.1/q, 0.5/0/o}
\tkzDefShiftPoint[o](60:0.6){g}
\tkzDefShiftPoint[o](60:3){a}
\fill[pattern=north east lines](m)--(n)--(p)--(q);
\tkzDrawSegments(m,n)
\tkzDrawSegments[gray, line width=3pt](a,o)
\tkzMarkAngles[size=0.4cm](n,o,a)
\tkzLabelAngles[pos=0.6](n,o,a){$\alpha$}
\draw [very thick,blue,->](a) -- ($(a)!1cm!-90:(o)$)node[above]{$\vec{F}$}coordinate(f);
\tkzMarkRightAngle(f,a,o)
\draw[very thick,blue,->](g)--++(-90:1.8)node[right]{$\vec{P}$};
\tkzDrawPoints[size=3,fill=white](o,a,g)
\tkzLabelPoint[below](o){O}
\tkzLabelPoint[right](a){A}
\tkzLabelPoint[above left](g){G}
\end{tikzpicture}}
} 
\end{ex}
\begin{ex}%book%Câu 6
\immini[thm]{Một tấm ván nặng $150\ N$ được bắc qua một con mương. Biết trọng tâm $G$ của tấm ván cách điểm tựa $A$ một khoảng là $2\ m$ và cách điểm tựa $B$ một khoảng $1\ m$. Lực mà tấm ván tác dụng lên hai bờ mương có độ lớn lần lượt là}{\begin{tikzpicture}[scale=1.3,thick,>=stealth',draw=Mapcolor]
\coordinate (a) at (0,0)node[above=6pt]{A};
\draw[] (a)rectangle(4,0.15);
\path (a)--++(0:4) coordinate(b)node[above=6pt]{B};
\draw (-0.5,0)-- (a)-- (0,-0.5);
\fill[pattern=north east lines] (-0.5,0)-- (a)-- (0,-0.5)-- (-0.1,-0.5)-- (-0.1,-0.1)-- (-0.5,-0.1);
\draw (4,-0.5)-- (b)-- (4.5,0);
\fill[pattern=north east lines] (4,-0.5)-- (b)-- (4.5,0)-- (4.5,-0.1)-- (4.1,-0.1)-- (4.1,-0.5);
\draw(2.5,0.07)circle(0.05cm)node[above=6pt]{G};
\end{tikzpicture}}
\choice
{\True $F_A=50\ N$; $F_B=100\ N$}
{$F_A=50\ N$; $F_B=150\ N$}
{$F_A=60\ N$; $F_B=100\ N$}
{$F_A=60\ N$; $F_B=150\ N$}
\haidong{10}	\loigiai{}
\end{ex}
\begin{ex}%book%[1P3B3-2][Loai1]
\immini[thm]{Một thanh nhẹ gắn vào sàn tại $B$. Tác dụng lên đầu $A$ lực kéo $ F=100\ N$ theo phương ngang. Thanh được giữ cân bằng nhờ dây $AC$.  Biết $\alpha=30^\circ$. Lực căng của dây có độ lớn là
\choice
{\True $200\ N$}
{$300\ N$}
{$400\ N$}
{$500\ N$}	}{\begin{tikzpicture}[scale=0.8,thick,>=stealth',draw=Mapcolor]
\coordinate (a) at (0,0) node[above]{\normalsize{A}};
\filldraw[black] (a) circle(1.2pt);  
\coordinate (b) at (0,-4) node[above]{\normalsize{}};
\filldraw[black] (b) circle(1.2pt)node[above right]{\normalsize{B}}; 
\draw[pattern=north east lines] (-0.1,0) rectangle (0.1,-4); 
\fill[pattern=north east lines] (-2.5,-4) rectangle (2.5,-4.2);
\draw (-2.5,-4)--(2.5,-4);
\draw[very thick] (a)--++(240:4.62) coordinate (c) node[above left]{\normalsize{C}};
\draw[very thick,->, Mapcolor] (a)--++(0:1.2) coordinate (t) node[above left]{\normalsize{\bth{\vec{F}}}};
\draw (a)++(240:0.6) arc (240:270:0.6) ;
\path (a) ++(252:0.9) node{\normalsize{\bth{\alpha}}};
\end{tikzpicture} }
\haidong{9}	
\loigiai{
\immini{Áp dụng quy tắc moment lực với trục quay qua B ta có:
$ M_F=M_T \Rightarrow F.AB=T.AB.sin\alpha \Rightarrow T=200\ N.$ }{\begin{tikzpicture}[scale=0.8,thick,>=stealth']
\coordinate (a) at (0,0) node[above]{\normalsize{A}};
\filldraw[black] (a) circle(1.2pt);  
\coordinate (b) at (0,-4) node[above]{\normalsize{}};
\filldraw[black] (b) circle(1.2pt)node[above right]{\normalsize{B}}; 
\draw[dashed] (b)--++(150:2) coordinate (d);
\draw (d)++(60:0.25)--++(330:0.25)--++(240:0.25);
\draw[pattern=north east lines] (-0.1,0) rectangle (0.1,-4); 
\fill[pattern=north east lines] (-2.5,-4) rectangle (2.5,-4.2);
\draw (-2.5,-4)--(2.5,-4);
\draw [very thick,->](0,-4)--(0,-1.825)node[right=5pt]{\normalsize{\bth{\vec{N}}}};
\draw[very thick] (a)--++(240:4.62) coordinate (c) node[above left]{\normalsize{C}};
\draw[very thick,->] (a)--++(240:2.4) coordinate (t) node[above left]{\normalsize{\bth{\vec{T}}}};
\draw[very thick,->] (a)--++(0:1.2) coordinate (t) node[above left]{\normalsize{\bth{\vec{F}}}};
\draw (a)++(240:0.6) arc (240:270:0.6) ;
\path (a) ++(252:0.9) node{\normalsize{\bth{\alpha}}};
\end{tikzpicture} }
}
\end{ex}













\subsubsection*{PHẦN 2. Thí sinh trả lời từ câu 1 đến câu 4. Trong mỗi ý a), b), c), d) ở mỗi câu, thí sinh chọn đúng hoặc sai.}
\setcounter{ex}{0}
\begin{ex}%Tailieuchuan
Một vật có khối lượng $5\ kg$ được buộc vào đầu một chiếc gậy dài $90\ cm$. Một người quẩy lên trên vai chiếc gậy đó biết vật cách vai một khoảng là $60\ cm$. Đầu còn lại của chiếc gậy được giữ bằng tay theo phương vuông góc với gậy. Bỏ qua trọng lượng của gậy, lấy $g=10 \ m/s^2$.
\choiceTF[t]
{\True Lực do vật tác dụng lên gậy là $50\ N$} 
{Moment lực do vật tác dụng lên gậy có độ lớn $300\ N.m$} 
{Để gậy được giữa cân bằng thì lực do vật tác dụng lên gậy phải bằng lực do tay tác dụng lên gậy} 
{Lực giữ của tay bằng $25\ N$} 
\loigiai{ 
\begin{itemchoice}
\itemch
\itemch
\itemch
\itemch
\end{itemchoice}
}
\end{ex}
\begin{ex}%Tailieuchuan
\immini[thm]{Hai chị em chơi bập bênh như hình bên dưới. Cho biết người chị có trọng lượng $P_2=400 \ N$, khoảng cách $d_2=1 \ m$, còn người em có trọng lượng $P_1=200 \ N$.}{
\includegraphics[width=4cm]{images/2025_04_24_a80008410be527bbd242g-11(1)}}
\choiceTF[t]
{\True Trọng lực của người chị làm bập bênh quay cùng chiều kim đồng hồ} 
{\True Vì bập bênh đang ở trạng thái cân bằng nên $P_2. d_2=P_1. d_1$} 
{Để bập bênh cân bằng như hình vẽ thì người em phải ngồi cách điểm quay $1,5\ m$} 
{Để bập bênh nghiêng về phía người em thì người em có thể ngồi cách người chị $2,8 \ m$} 
\loigiai{ 
\begin{itemchoice}
\itemch
\itemch
\itemch
\itemch
\end{itemchoice}
}
\end{ex}
\begin{ex}%Tailieuchuan
\immini[thm]{Một người nâng tấm ván $AB$ có khối lượng $40\ kg$ (phân bố đều theo chiều dài) với lực $F$ để ván nằm yên và hợp với mặt đường một góc $30^\circ$. Lấy $g=10 \ m/s^2$ và hướng của lực $F$ vuông góc với tấm ván.}{
\begin{tikzpicture}[scale=1,thick,>=stealth',draw=Mapcolor]
\tkzInit[ymin=-0.5,ymax=4,xmin=-0,xmax=4]\tkzClip
\tkzDefPoints{0/0/m, 3.5/0/n, 3.5/-0.1/p, 0/-0.1/q, 0.5/0/o}
\tkzDefShiftPoint[o](60:0.6){g}
\tkzDefShiftPoint[o](30:3){a}
\fill[pattern=north east lines](m)--(n)--(p)--(q);
\tkzDrawSegments(m,n)
\tkzDrawSegments[gray, line width=3pt](a,o)
\tkzMarkAngles[size=0.7cm](n,o,a)
\tkzLabelAngles[pos=1.0](n,o,a){$30^\circ$}
\draw [very thick,Mapcolor,->](a) -- ($(a)!1cm!-90:(o)$)node[above]{$\vec{F}$}coordinate(f);
\tkzDrawPoints[size=3,fill=white](o,a)
\tkzLabelPoint[above left](o){$A$}
\tkzLabelPoint[right](a){$B$}
\draw[very thick,->,Mapcolor] (2,0.9) -- (2,0.2);
\node at (2.5,0.5) {$\vec{P}$};
\node at (1.2,-0.3) {$d_P$};
\end{tikzpicture}}
\choiceTF[t]
{\True Moment lực $F$ có tác dụng làm ván quay ngược chiều kim đồng hồ} 
{\True Moment lực $P$ có tác dụng làm ván quay cùng chiều kim đồng hồ} 
{\True Trọng lực $P$ có điểm đặt ở chính giữa tấm ván} 
{\True Lực $F$ có độ lớn là $100 \sqrt{3}\ N$} 
\loigiai{ 
\begin{itemchoice}
\itemch
\itemch
\itemch
\itemch
\end{itemchoice}
}
\end{ex}
\begin{ex}%Tailieuchuan
Xét một lực tác dụng lên vật có trục quay cố định.
\choiceTF[t]
{\True Moment của lực đối với một trục quay phụ thuộc vào độ lớn của lực và khoảng cách từ điểm tác dụng của lực đến trục quay} 
{Moment của lực là đại lượng vô hướng và chỉ phụ thuộc vào độ lớn của lực} 
{\True Nếu lực đi qua trục quay, moment của lực bằng $0$} 
{Lực có độ lớn $5\ N$ thì moment bằng $5 \ N. m$} 
\loigiai{ 
\begin{itemchoice}
\itemch
\itemch
\itemch
\itemch
\end{itemchoice}
}
\end{ex}





















\subsubsection*{PHẦN 3. Thí sinh trả lời từ câu 1 đến câu 6.}
\setcounter{ex}{0} 
\begin{ex}%Tailieuchuan
Để có moment của một vật có trục quay cố định là $10\ N.m$ thì cần phải tác dụng vào vật một lực bằng bao nhiêu niutơn biết khoảng cách từ giá của lực đến tâm quay là $20\ cm$?
\shortans[oly]{50}
\loigiai{}
\end{ex}
\begin{ex}%Tailieuchuan
\immini[thm]{Thước dẹt, đồng chất, tiết diện đều, trọng lượng $20\ N$ và quay quanh trục $O$. Biết $OG=40\ cm$ và thước hợp với đường thẳng đứng qua $O$ một góc $45^\circ$. Moment trọng lực của thước là bao nhiêu $N.m$?}{
\begin{tikzpicture}[scale=1,thick,>=stealth',draw=Mapcolor]
\path
(0,-0.4)coordinate(O)
(1.1,-1.5)coordinate(G)
;
\begin{scope}[shift={(0,0)},rotate=-45]
\draw[fill=Mapcolor!30]  (0,0) rectangle (3.5,-0.5);
\end{scope}

\draw[dashed] (O) -- (0,-2.9);

\foreach \x/\g in {O/90,G/90}
\draw[fill=red](\x) circle(0.075)+(\g:0.3)node{$\x$};
\end{tikzpicture}}
\shortans[oly]{5,66}
\loigiai{
}
\end{ex}
\begin{ex}%Tailieuchuan
Một lực có phương song song với trục quay của vật rắn và có độ lớn là $5,5\ N$. Biết khoảng cách từ giá của lực tới trục quay là $2\ m$. Moment lực của một lực đối với trục quay bằng bao nhiêu niutơn?
\shortans[oly]{0}
\loigiai{}
\end{ex}
\begin{ex}%Tailieuchuan
\immini[thm]{Cho đòn bẩy như hình vẽ bên dưới. Điểm đặt trọng lực của đòn bẩy vào đúng điểm $O$. Đầu $A$ của đòn bẩy treo một vật có trọng lượng $30\ N$. Chiều dài đòn bẩy là $50\ cm$. Khoảng cách từ đầu $A$ đến trục quay là $20\ cm$. Đầu $B$ của đòn bẩy phải treo một vật khác có trọng lượng bằng bao nhiêu niutơn để đòn bẩy cân bằng như ban đầu?}{
\begin{tikzpicture}[scale=1,draw=Mapcolor, very thick,>=stealth',draw=Mapcolor]
\tkzDefPoints{0/0/o1, 5/0/o2, 3.475/0.15/c, 2.25/0.15/g,0/-1.5/m, 3.475/-1.25/m'}
\draw(-0.05,0.3)--(-0.05,0)--(4,0)--(4,0.3)--cycle;

\node at (m) [shape=rectangle,draw=red,pattern=dots,minimum height=0.5cm,minimum width=0.7cm, rounded corners=1pt,label=left:] {};

\node at (-0.1,0.5) {A};
\node at (1,0.5) {O};
\draw (0,0) -- (0,-1.22);
\node at (4,0.5) {B};
\draw[fill=gray!30] (1,0) -- (0.9,-0.2) -- (1.1,-0.2) -- (1,0);
\end{tikzpicture}}
\shortans[oly]{20}
\loigiai{
}
\end{ex}
\begin{ex}%Tailieuchuan
Một vật rắn chịu tác dụng của lực $F=20 \ N$ có thể quay quanh trục cố định, khoảng cách từ giá của lực đến trục quay là $20\ cm$. Moment của lực $F$ tác dụng lên vật là bao nhiêu $N.m$?
\shortans[oly]{4}
\loigiai{}
\end{ex}
\begin{ex}%Tailieuchuan
Một thước thẳng, đồng chất, tiết diện đều có chiều dài $AB=100 \ cm$, trọng lượng $P=30 \ N$. Thước có thể quay xung quanh một trục nằm ngang đi qua điểm $O$ trên thước với $OA=30 \ cm$. Để thước cân bằng nằm ngang, cần treo tại đầu $A$ một vật có trọng lượng bằng bao nhiêu niutơn?
\shortans[oly]{20}
\loigiai{}
\end{ex}
















